\begin{helper}
In a fixed history cell, assume each complete blue terminal branch supplies a
red exceptional set \(E_R(\beta)\) satisfying the closed-pair bound
\(|E_R(\beta)|\le 3\varepsilon_1 k^2\) and covering every red present
staged-open coordinate of any regular independent completion of that branch.
Then there is one natural number \(N_R\), independent of the branch, such that
\[
N_R\le 3\varepsilon_1 k^2,
\]
and every nonempty blue branch admits a red exceptional set of cardinality at
most \(N_R\) with the same red coverage property.
\end{helper}
\begin{proof}
The sample space is finite, because a two-bite sample is a pair of finite
graphs together with a finite injection.  Hence the set of history branches is
finite, and so is the subfamily of branches that actually occur in the
fixed-count argument.

For each branch in this finite family, choose one red exceptional set supplied
by the branch hypothesis.  The branch hypothesis gives two facts for that
chosen set: it is contained in the terminal coordinate set, and its
cardinality satisfies
\[
|E_R(\beta)|\le 3\varepsilon_1 k^2
\]
by the definition of \(\noderef{ClosedPairsBound}\).  It also gives the
coverage statement: any regular completion which agrees with the branch on
the embedding, on recorded coordinates, and on all blue terminal coordinates
has all red present staged-open coordinates inside \(E_R(\beta)\).

Let \(N_R\) be the maximum of the finitely many cardinalities
\(|E_R(\beta)|\), and take \(N_R=0\) if there are no relevant branches.  Since
each member of the finite set being maximized is at most
\(3\varepsilon_1 k^2\), the maximum is also at most
\(3\varepsilon_1 k^2\).  For a given branch, its chosen exceptional set has
cardinality at most this maximum and retains the terminal-containment and
red-coverage properties from the original branch hypothesis.  This is the
desired global red witness.
\end{proof}
